\section{Residually reducible comparison theorems: exact applicability}
\label{sec:kt-reducible-comparison}

This audit concerns the integral extension module at the residual
point $\bar\rho=\bar\chi^e\oplus1\oplus\bar\chi^{-e}$.
Here $e>0$ is even, $\chi=\varepsilon_\ell$, $\ell>3$, and the
carried GD condition gives $(\ell-1)\nmid2e$.  The desired module
is the odd weight $-1$ cotangent identified with the original
rational strict Selmer dual in KT-ODD-CONORMAL-SELMER-EQUIVALENCE.
The conclusions below are an applicability audit, not a new
upper bound for that Selmer group.

\subsection{What the general matrix algebra results actually control}

Bella\"iche--Chenevier, \emph{$p$-adic families of Galois
representations and higher rank Selmer groups}, arXiv
\texttt{math/0602340v2}, 12 January 2007, is the exact author
version used here.  Its pagination is not that of the final
Ast\'erisque volume.  Proposition~1.5.1, pp.~32--34, applies
to any residually multiplicity-free pseudocharacter.  Thus
three distinct rank-one constituents cause no obstruction.
For a Cayley--Hamilton quotient with structural modules $A_{ij}$,
put
\[
 A'_{ij}=\sum_{k\ne i,j}A_{ik}A_{kj},\qquad
 I=\sum_{i\ne j}T(A_{ij}A_{ji}).                         \tag{RC.1}
\]
The total reducibility ideal $I$ is independent of that
Cayley--Hamilton quotient.  Theorem~1.5.5, p.~35, gives
\[
 \operatorname{Hom}_A(A_{ij}/A'_{ij},A/J)
       \lhook\joinrel\longrightarrow
       \operatorname{Ext}^1_{R/JR}(\rho_j,\rho_i),
       \qquad J\supset I.                                \tag{RC.2}
\]
Theorem~1.5.6, p.~36, identifies the image with extensions
factoring through the chosen Cayley--Hamilton quotient.
Remark~1.5.8, p.~38, explicitly distinguishes these extension
groups from those of the faithful trace quotient: their
inclusion can be strict even though the reducibility ideal
is unchanged.  Consequently a comparison of trace rings
alone supplies neither an identification of the required
off-diagonal module nor its integral saturation.

\begin{lemma}[KT-GMA-REDUCIBILITY-DOES-NOT-BOUND-EXTENSIONS]
For every integer $n\geq1$ there is a polarized
Cayley--Hamilton generalized matrix algebra of type $(1,1,1)$
over a complete discrete valuation ring $O$, with base
algebra $O$ and total reducibility ideal zero, whose two
adjacent off-diagonal modules have length $n$.
\end{lemma}
\begin{proof}
Put $T_n=O/(\varpi^n)$ and
$R_n=O\oplus\epsilon T_n$, with $\epsilon^2=0$.
Inside $M_3(R_n)$ take
\[
 \mathcal E_n=
 \left\{
 \begin{pmatrix}
 a&0&0\\ \epsilon u&b&0\\0&\epsilon v&c
 \end{pmatrix}:
 a,b,c\in O,\ u,v\in T_n\right\}.                       \tag{RC.3}
\]
It is closed under multiplication; every product of two
off-diagonal structural modules is zero.  Diagonal
evaluation defines a degree-three determinant law, with
characteristic polynomial $(X-a)(X-b)(X-c)$.
Cayley--Hamilton holds in the displayed matrix algebra,
also after every scalar extension, since the same direct
sum description embeds in the corresponding matrix
algebra after base change.  The anti-automorphism
$x\mapsto Jx^tJ$, with $J$ anti-diagonal, exchanges
$e_1,e_3$ and $A_{21},A_{32}$ and fixes $e_2$.
Thus the polarization has the required outer exchange.
Formula (RC.1) gives $I=0$, while
$A_{21}/A'_{21}\simeq A_{32}/A'_{32}\simeq T_n$.
The coefficient algebra $O$ is itself flat and a complete
intersection.  This is an algebraic scope example; it is
not asserted to satisfy the global automorphic source
conditions.
\end{proof}

\subsection{The earlier two-constituent comparisons}

Berger--Klosin, \emph{On deformation rings of residually
reducible Galois representations and $R=T$ theorems},
author version \texttt{hd32.pdf}, dated 10 February 2011,
uses two nonisomorphic absolutely irreducible residual
blocks and a nonsplit residual extension.  Assumption~5.1,
p.~13, imposes a one-dimensional extension Selmer group
and rigid diagonal deformation rings.  Theorem~8.5,
p.~23, requires an upper bound for the divisible extension
Selmer group in addition to a lower bound for the Hecke
congruence quotient.  The alternative Theorem~8.6 assumes
the needed conormal inequality.  Grouping two characters
of our residual triple into one block does not satisfy
absolute irreducibility.  Even for a two-character
subproblem these hypotheses do not supply the missing
upper bound: that bound is an input.

Their \emph{On modularity of reducible residual Galois
representations}, author version
\texttt{bergerklosinonlifting21.pdf}, dated 2 April 2013,
retains two absolutely irreducible blocks.  Assumption~2.7,
p.~6, is explicitly the required Bloch--Kato upper bound;
Assumption~4.2 and Lemma~4.4, p.~12, use it together with
a congruence lower bound to prove surjectivity onto the
Selmer group.  It cannot be removed when using that
surjectivity.  The construction also passes to the image
of the reduced ring in its localization at $p$; the
introduction, pp.~1--3, explains why characteristic-$p$
components can be lost.  This does not identify our full
integral cotangent before that passage.

Urban's \emph{Selmer Groups and the Eisenstein--Klingen
Ideal}, arXiv \texttt{math/9803171v1}, Theorem~1.1,
author p.~6, does provide a useful two-block stable
lattice.  Its nonsplit residual extension yields a
Selmer class.  This is a construction of classes, not
an upper bound for every class.  We use precisely this
positive consequence below.

Allen--Newton--Thorne, \emph{Automorphy lifting for
residually reducible $l$-adic Galois representations, II},
Theorem~1.1 (also Theorem~6.1), allows more constituents,
but its conditions (iv) and (vii) fail here.  After the
rank-three normalization by $\bar\chi^{-1}$, individual
polarization of either outer constituent requires
$\bar\chi^{2e}=1$, contrary to GD.  All three characters
become equal on $G_{E(\zeta_\ell)}$, contrary to (vii).
The full triple's polarization exchanging outer
constituents does not repair either condition.  The
independent source audit is recorded in
\texttt{lean\_chain/thorne\_scope/THORNE\_SCOPE.md}.

\subsection{The three-character ordinary unitary preprint}

Berger--Klosin, \emph{On $p$-adic modularity of abelian
threefolds with imaginary multiplication and rational
$p$-isogeny}, arXiv \texttt{2611.19007v1}, 13 November
2026, is a recent preprint, not a claimed published
theorem.  The official arXiv PDF and the separately
retrieved mirror are byte-identical.  Its general
Proposition~2.12 and Corollary~2.16, pp.~7--9, bound
$I/I^2$ using off-diagonal modules under vanishing and
cyclicity hypotheses.  The direction of this estimate
does not bound those modules from $I/I^2$.

Section~10, pp.~51--56, is directly relevant but has
different ordinary data.  Write $\Psi$ for its
specialized character.  Its outer characters are
$\tau_1=\Psi\chi^{-1}$ and
$\tau_3=\Psi^{-c}\chi$, with middle constituent $1$.
Matching our characteristic-zero triple gives just
the following two possibilities:
\[
 \begin{array}{c|c|c}
 \tau_1&\Psi\text{ infinity type }(a+b,-b)&(a,b)\\ \hline
 \chi^e&(e+1,e+1)&(2e+2,-e-1)\\
 \chi^{-e}&(1-e,1-e)&(2-2e,e-1).
 \end{array}                                                \tag{RC.4}
\]
The first violates $b\geq-1$, the second $a\geq3$.
These are exact characteristic-zero weights, so reducing
the exponents modulo $\ell-1$ does not resolve the issue.

There is a separate local flag mismatch.  Equations
(10.1)--(10.2), p.~52, prescribe the upper-triangular
orders $(\tau_1,\tau_3,1)$ and $(1,\tau_1,\tau_3)$ at
the two split primes.  Our component has order
$(\chi^e,1,\chi^{-e})$ at both.  An identification
respecting polarization fixes the unique self-dual
constituent $1$ and can at most exchange the outer two;
it cannot identify these ordered functors.  A locally
split representation can have several flags, so this
observation does not assert that every individual point
fails the preprint's flag condition.  Our construction
does not supply that additional pair of flags.

For clarity, the following substitution is a
\emph{hypothesis audit outside the theorem's weight
range}, not an application.  Substituting
$\Psi=\chi^{e+1}$ into the Selmer conditions of
Theorem~10.2, p.~53, gives
\[
 \begin{array}{c|c}
 \text{condition}&\text{coefficient character}\\ \hline
 (S2)(a)&\chi^e\\
 (S2)(b),\ c=-1&\eta\chi^{2e}\text{ over }\mathbf Q\\
 (S2)(c)&\chi^{-e}\\
 (S3)&(K_O/O)(\chi^{-e}).
 \end{array}                                                \tag{RC.5}
\]
The class component in Theorem~10.4(5) similarly has
character $\eta\bar\chi^{2e}$; it is not literally the
untwisted target.  Condition~10.4(6), however, retains
the same one-dimensional $\chi^{-e}$ Selmer assumption.
Assumption (T), p.~53, independently requires reduced
specializations of the relevant ordinary Hecke family.
No such family comparison or reducedness hypothesis has
been established for our ring.

Definition~8.22, p.~48, places inertia conditions at
\emph{both} split primes on a parallel negative character
$\chi^{-e}$.  The one-sided examples in that section
belong to the preprint's different infinity-type cone.
In either interpretation the target strict group injects
into the group in (RC.5).  The proof of Theorem~10.2,
pp.~53--55, uses its Selmer upper bounds to bound
$I/I^2$ and then applies a numerical comparison.
It does not extract those upper bounds from $R=T$.
The fixed prime $697446897311$ in the $e=26$ calculation
is congruent to $2$ modulo $3$, hence is inert in
$\mathbf Q(\sqrt{-3})$; the preprint's split-prime
assumption already excludes that particular seed.

\subsection{The cyclicity condition already detects the target}

\begin{theorem}[KT-RESIDUAL-CYCLICITY-SCOPE]
Let $E=\mathbf Q(\sqrt{-3})$, $\eta=\chi_{-3}$,
and assume that $\ell$ splits in $E$.  Let $f$ be one
of the non-CM ordinary forms supplied by the checked
lower-weight Eisenstein congruence: its two-dimensional
representation $\rho_f$ is absolutely irreducible,
\[
 \bar\rho_f^{\rm ss}=1\oplus\eta\bar\chi^e,
 \qquad a_\ell(f)\equiv1,
                                                               \tag{RC.6}
\]
and $\rho_f|_{G_E}$ is unramified away from $\ell$.
Let $k$ be its residue field.  Let $\Sigma$ be any
conjugation-stable set containing the primes over
$3\ell$, and set
\[
 \mathcal H_E=
 \ker\left(H^1(G_{E,\Sigma},k(\bar\chi^{-e}))
       \longrightarrow\prod_{v\mid\ell}
                         H^1(I_v,k(\bar\chi^{-e}))\right).
                                                               \tag{RC.7}
\]
Let $\mathcal X$ be the original rational strict
Selmer group with coefficients
$(\mathbf Q_\ell/\mathbf Z_\ell)(\chi^{-e})$,
including the decomposition-group condition at $2$.
Then
\[
 \dim_k\mathcal H_E^{c=-1}\geq1,\qquad
 \dim_k\mathcal H_E^{c=+1}
          \geq\dim_{\mathbf F_\ell}\mathcal X[\ell].
                                                               \tag{RC.8}
\]
In particular, $\dim_k\mathcal H_E=1$ implies
$\mathcal X=0$.  The same implication holds with any
relaxation of the local conditions in (RC.7).
Thus an attempted use of the cyclicity hypothesis in
(RC.5) in these concrete branches would already require
a condition forcing the original vanishing problem.
\end{theorem}

\begin{proof}
Apply Urban's Theorem~1.1 with coefficient ring the
valuation ring of $f$, ideal its maximal ideal, and
ordered rank-one blocks $1$, $\eta\chi^e$.
Integrality and congruence of characteristic polynomials
follow from (RC.6).  Absolute irreducibility and
distinctness of the residual blocks are part of the
hypotheses, the latter also following from GD at
inertia over $\ell$.  The theorem supplies a stable
lattice whose reduction is a nonsplit extension
\[
 0\longrightarrow k\longrightarrow\bar L
       \longrightarrow k(\eta\bar\chi^e)
       \longrightarrow0.                                  \tag{RC.9}
\]
Urban's auxiliary module is a nonzero finitely generated
submodule of the coefficient field, hence is free of
rank one over the valuation ring; its reduction is
exactly one copy of $k$ as in (RC.9).
Indeed its assertion that the lattice has no quotient
isomorphic to the first residual block rules out a
split reduction.  This gives a nonzero class
$u\in H^1(G_{\mathbf Q},k(\eta\bar\chi^{-e}))$.

Ordinarity gives a characteristic-zero local stable
line whose character reduces to $\eta\bar\chi^e$:
the ordinary quotient reduces to $1$ because
$a_\ell(f)\equiv1$.  Its intersection with any stable
lattice is saturated.  In (RC.9) its reduction and the
displayed stable line $k$ are distinct, since their
inertial characters differ.  They therefore split
(RC.9) on the full decomposition group at $\ell$.
Thus $u$ is locally strict there, not merely
unramified.  Away from $3\ell$ it is unramified.
At three, restriction to $G_E$ is unramified by the
stated property of $\rho_f$, independently of lattice.
Consequently $u|_{G_E}$ belongs to (RC.7).

Restriction through the index-two subgroup $G_E$ is
injective because two is a unit.  Under the
identification of the restricted coefficient character
with $\bar\chi^{-e}$, the additional factor $\eta$
places its image in the $c=-1$ eigenspace.  Hence the
first inequality in (RC.8).

Put $M=(\mathbf Q_\ell/\mathbf Z_\ell)(\chi^{-e})$.
The global and local $H^0$ groups at $\ell$ vanish:
GD supplies an inertia element on which
$\bar\chi^{-e}-1$ is a unit.  The exact sequence
$0\to\mathbf F_\ell(\bar\chi^{-e})\to M
\xrightarrow{\ell}M\to0$ therefore identifies
global $\ell$-torsion classes with finite-coefficient
classes and preserves their strict local vanishing
at $\ell$.  Away from $\ell$, the action on inertia
is trivial, so $H^0(I_q,M)=M$ is $\ell$-divisible.
The same long exact sequence therefore makes
$H^1(I_q,\mathbf F_\ell(\bar\chi^{-e}))\to H^1(I_q,M)$
injective, preserving unramifiedness.  Restricting these finite classes to
$E$ and extending scalars to $k$ gives an injection
\[
 \mathcal X[\ell]\otimes_{\mathbf F_\ell}k
                  \lhook\joinrel\longrightarrow
                         \mathcal H_E^{c=+1}.              \tag{RC.10}
\]
No vanishing of finite-coefficient unramified cohomology
at $2$ is asserted: the extra condition at $2$ is
simply part of the domain $\mathcal X$.
This proves the second inequality.  Since the two
eigenspaces are disjoint and $\ell\ne2$, dimension
one for (RC.7) forces $\mathcal X[\ell]=0$.
Every nonzero element of the $\ell$-primary torsion
group $\mathcal X$ has a nonzero multiple of order
$\ell$, so $\mathcal X=0$.
\end{proof}

The theorem uses the already checked existence,
ordinarity and irreducibility of the classical form;
it does not supply the cyclicity assumption.
The inspected comparison results therefore leave a
specific missing input: an integral comparison of
the relevant extension module, or an independent
upper bound for it, with the actual ordinary flags
and conjugation sign.  Neither a semisimple $R=T$
statement nor a reducibility-ideal estimate reverses
the injections in (RC.2).
